% =========================================================
% Order-Isomorphisms
% =========================================================

\section{Order-Isomorphisms}

In the study of partially ordered sets (posets), it is essential to identify when two ordered sets $(P, \le_P)$ and $(Q, \le_Q)$ have the same order-theoretic structure. We say that $P$ and $Q$ are \textbf{order-isomorphic}, and write $P \cong Q$, if there exists a map $\varphi: P \to Q$ that preserves and reflects the order.

\begin{definition}[Order-Isomorphism]
A map $\varphi: P \to Q$ is an \textbf{order-isomorphism} if it is a bijection and satisfies, for all $x, y \in P$,
\[
x \le_P y \iff \varphi(x) \le_Q \varphi(y).
\]
Such a map preserves the order structure exactly.
\end{definition}

\subsection*{Consequences and Properties}

\begin{enumerate}
    \item \textbf{Injectivity from Order Reflection:}  
    If $\varphi: P \to Q$ satisfies
    \[
    x \le_P y \iff \varphi(x) \le_Q \varphi(y)
    \]
    and is surjective, then it is automatically injective. Indeed, if $\varphi(x) = \varphi(y)$, then
    \[
    \varphi(x) \le_Q \varphi(y) \quad \text{and} \quad \varphi(y) \le_Q \varphi(x),
    \]
    hence
    \[
    x \le_P y \quad \text{and} \quad y \le_P x,
    \]
    and by antisymmetry, $x = y$.

    \item \textbf{Stability under Inversion:}  
    If $\varphi: P \to Q$ is an order-isomorphism, then its inverse $\varphi^{-1}: Q \to P$ is also an order-isomorphism.
\end{enumerate}

\begin{remark}[Isomorphism vs.\ Bijection]
A bijection between posets need not be an order-isomorphism. For example, let $P = \{0,1\}$ with $0 < 1$ and $Q = \{0,1\}$ with $0 < 1$. The bijection defined by $f(0)=1$ and $f(1)=0$ is not an order-isomorphism, since it reverses the order: $0 < 1$ but $f(0) \not\le f(1)$.
\end{remark}